\chapter{相关理论与关键技术}
\label{cha:theory}

本章围绕本文方法所涉及的核心理论与关键技术展开介绍。首先，对 POI 序列预测任务的基本概念、任务特点及建模难点进行说明，明确本文研究问题的理论背景。随后，介绍残差量化变分自编码器（Residual Quantized Variational AutoEncoder, RQ-VAE）的基本原理，说明其在离散语义表示学习中的作用。接着，对大语言模型及参数高效微调技术进行概述，为后续基于生成式框架的 POI 序列预测提供技术基础。最后，介绍地理空间散列编码方法，重点说明其在位置离散表示和空间邻近性建模中的应用。通过本章的论述，可以为后续章节中的方法设计与实验分析奠定理论基础。

\section{POI 序列预测任务概述}
\label{sec:poi_sequence_prediction}

兴趣点（Point-of-Interest, POI）通常指地图服务、社交签到系统或本地生活平台中的特定地点实体，例如餐厅、商场、车站、学校和景点等。在基于位置的服务场景中，用户的历史行为往往表现为一系列按时间顺序排列的 POI 访问记录。这些记录既反映了用户的兴趣偏好，又包含了明显的时空约束信息，因此基于用户历史访问序列预测其未来可能访问的地点，构成了 POI 序列预测任务的核心问题。

从形式化角度来看，设用户集合为 $\mathcal{U}=\{u_1,u_2,\dots,u_M\}$，POI 集合为 $\mathcal{P}=\{p_1,p_2,\dots,p_N\}$。对于任意用户 $u$，其历史访问序列可以表示为
\begin{equation}
S_u=\{(p_1,t_1),(p_2,t_2),\dots,(p_T,t_T)\},
\end{equation}
其中，$p_i \in \mathcal{P}$ 表示第 $i$ 次访问的兴趣点，$t_i$ 表示对应的访问时间，$T$ 为该用户历史访问序列长度。POI 序列预测的目标是在给定用户历史访问序列以及目标时刻上下文的条件下，预测用户下一时刻最可能访问的目标兴趣点 $p_{T+1}$。因此，该任务本质上属于一种具有时空约束的序列推荐问题。

与一般序列推荐任务相比，POI 序列预测具有更为复杂的特性。首先，用户访问地点的行为不仅由其长期兴趣决定，还会受到强烈的空间约束。例如，用户短时间内连续访问两个地理距离极远的地点在现实中往往不具备可行性，因此 POI 之间的地理邻近性对预测结果具有重要影响。其次，用户访问行为通常具有显著的时间规律。工作日与周末、白天与夜晚、通勤时段与休闲时段下，用户访问 POI 的偏好可能存在明显差异。再次，POI 数据天然具有稀疏性和多样性。一方面，单个用户的访问序列通常较短，重复访问频率有限；另一方面，城市中的 POI 数量巨大，不同 POI 之间的类别差异、地理分布和语义属性都较为复杂。这些因素共同增加了序列预测任务的难度。

围绕上述问题，研究者提出了多种 Next POI 推荐方法。早期工作多基于马尔可夫链、矩阵分解和时空概率模型，通过显式建模地点转移概率、距离衰减和用户偏好完成预测。这类方法具有较好的可解释性，但在面对复杂时空依赖关系时表达能力有限。随着深度学习的发展，研究者开始采用循环神经网络、注意力机制和 Transformer 等结构建模用户访问序列。此类方法在捕捉长短期兴趣转移、时间依赖以及上下文交互方面表现出更强能力，并推动了 POI 序列预测研究的快速发展 \cite{feng2018deepmove,yang2020flashback,luo2021stan,kang2018sasrec}。

尽管如此，现有方法仍然面临两个突出问题。其一，大多数模型主要关注用户历史访问顺序及上下文建模，而对 POI 自身表示的研究相对不足。传统做法通常将 POI 视作离散 ID，或简单利用类别、区域等文本属性进行编码，这种方式往往难以充分反映不同 POI 之间的语义联系。其二，地理信息虽然被普遍认为是 POI 推荐中的关键因素，但其引入方式仍较为粗糙，常见做法包括将经纬度直接作为数值特征输入模型，或将距离作为额外约束项加入损失函数。这些方式虽然能够在一定程度上引入空间信息，但难以与 POI 的语义表示形成统一表示空间，从而影响模型对语义--空间联合结构的建模能力。

近年来，随着生成式推荐和离散语义表示的发展，研究者开始尝试将推荐任务转化为序列生成问题，通过生成式模型直接生成目标物品的离散标识。在这一背景下，如何为 POI 构造更合理的离散表示，成为影响模型性能的关键因素。GNPR-SID 等方法的提出表明，相较于随机数字 ID，具有层级结构和语义关联的离散语义标识（Semantic ID, SID）更适合生成式推荐模型建模 \cite{rajput2023tiger,wang2025gnprsid}。然而，对于 POI 这种天然具有强空间属性的对象，如果离散语义表示仅由类别名称或区域文本学习得到，那么得到的 SID 更多反映的是语义属性，而无法充分刻画空间差异。也正因如此，在 POI 序列预测任务中，将地理信息融入离散表示学习过程具有较强的研究价值。

本文的工作正是在这一背景下展开。相较于将地理标识直接作为显式输出目标进行联合生成，本文更关注在表示学习阶段引入地理信息增强表示，将包含空间属性的 $\mathrm{geo\_emb}$ 输入 RQ-VAE，以学习更具空间感知能力的 SID，从而为后续基于 GNPR-SID 的生成式 POI 序列预测提供更有效的目标表示。

\section{RQ-VAE 原理}
\label{sec:rqvae}

变分自编码器（Variational AutoEncoder, VAE）是一类经典的生成模型，其核心思想是通过编码器将输入样本映射到潜在空间，再利用解码器从潜在表示中恢复原始输入，从而学习数据的潜在分布结构。在推荐系统和表示学习任务中，VAE 被广泛用于连续潜变量建模。然而，生成式推荐场景下的目标通常是离散 token 序列，因此如何将连续表示映射为具有语义结构的离散标识，成为一个关键问题。向量量化（Vector Quantization, VQ）及其扩展模型正是在这一需求下得到关注。

VQ-VAE 通过引入离散码本（codebook），将编码器输出映射为码本中最接近的离散码向量，从而实现“连续输入--离散表示”的转换。对于输入样本 $x$，编码器首先得到连续潜在表示 $z_e(x)$，然后在码本 $\mathcal{E}=\{e_1,e_2,\dots,e_K\}$ 中选择与其距离最近的码向量：
\begin{equation}
z_q(x)=e_k,\quad k=\arg\min_j \|z_e(x)-e_j\|_2.
\end{equation}
量化后的表示 $z_q(x)$ 会被送入解码器进行重构。通过同时优化重构损失与码本学习目标，模型可以逐步形成稳定的离散表示结构。

尽管 VQ-VAE 已经能够实现离散表示学习，但其单级码本通常受限于表达容量。当目标对象数量较大、语义结构复杂时，仅使用单个码本向量往往难以充分表示细粒度差异。为解决这一问题，研究者提出了残差量化变分自编码器（Residual Quantized Variational AutoEncoder, RQ-VAE）。该方法借鉴残差量化思想，不再依赖单一量化结果，而是通过多级码本逐层逼近输入表示。其基本过程可以描述如下：对于输入表示 $x$，编码器得到连续潜表示 $z_e(x)$，第一级码本选择最接近的量化向量 $e_{k_1}^{(1)}$，随后计算残差
\begin{equation}
r_1=z_e(x)-e_{k_1}^{(1)}.
\end{equation}
第二级码本再对残差 $r_1$ 进行量化，得到
\begin{equation}
e_{k_2}^{(2)}=\arg\min_j \|r_1-e_j^{(2)}\|_2.
\end{equation}
重复这一过程，经过 $M$ 级量化后，最终量化表示可以写为
\begin{equation}
z_q(x)=\sum_{m=1}^{M} e_{k_m}^{(m)}.
\end{equation}
对应的离散索引序列
\begin{equation}
(k_1,k_2,\dots,k_M)
\end{equation}
便可作为输入对象的多层级离散表示。

与单级 VQ-VAE 相比，RQ-VAE 的最大特点在于其离散表示具有更强的组合表达能力。由于每一级码本负责逼近上一层残差，量化过程天然形成从粗到细的层级结构。对于推荐系统中的物品表示学习而言，这意味着模型能够利用多级离散编码更精细地表达物品间的语义差异。特别是在语义 ID 学习场景中，RQ-VAE 生成的索引序列可以直接作为物品的语义 token 表示，这为后续的生成式模型训练提供了天然的离散目标形式。

在 GNPR-SID 这类生成式 POI 推荐框架中，RQ-VAE 的作用可以概括为：将原始的 POI 连续表示压缩并离散化为多层级语义标识 SID。设某一 POI 的输入表示为 $x_p$，经过 RQ-VAE 后得到的离散索引序列记为
\begin{equation}
SID(p)=\{s_1,s_2,\dots,s_M\},
\end{equation}
其中每个 $s_i$ 对应第 $i$ 级码本选出的索引。这样一来，原本难以直接作为生成目标的 POI 连续表示，便可以转化为适合大语言模型建模的离散 token 序列。

RQ-VAE 的另一个重要优势在于，其学得的多级离散表示通常能够在一定程度上反映样本间的层级结构。语义相近或分布相近的对象，往往在较高层级码本中共享更多相同或相似的离散索引，而更深层级则进一步区分细粒度差异。这种“由粗到细”的表达方式，对于生成式推荐任务具有重要意义：一方面，它降低了模型直接生成高维连续向量或随机物品 ID 的难度；另一方面，它使目标表示在 token 空间中具备一定的结构性，从而有助于模型学习上下文与目标之间更稳定的映射关系。

对于本文任务而言，RQ-VAE 之所以关键，在于它不仅承担离散化编码的角色，还连接了“表示学习”与“生成式预测”两个阶段。若输入给 RQ-VAE 的仅仅是传统的类别名称或区域文本所对应的语义表示，则最终得到的 SID 主要反映的是类别层面的语义信息。为了提升 SID 对 POI 空间属性的感知能力，本文将地理信息增强表示 $\mathrm{geo\_emb}$ 作为 RQ-VAE 的输入，使离散语义标识的学习从源头上融入空间信息。这样得到的 SID 便不再只是纯语义编码，而是具有一定地理区分能力的离散语义表示，这也是本文方法设计的核心出发点之一。

\section{大语言模型与参数高效微调}
\label{sec:llm_peft}

大语言模型（Large Language Models, LLMs）是近年来自然语言处理和推荐系统研究中的重要进展。此类模型通常基于 Transformer 架构，通过在大规模语料上进行自监督预训练，学习通用的语言理解与生成能力。与传统的任务专用模型相比，大语言模型具有更强的上下文建模能力、表达能力和迁移能力，因此逐渐被引入推荐系统、检索系统以及用户行为建模等任务中 \cite{brown2020gpt3,vaswani2017attention}。

从结构上看，Transformer 的核心机制是自注意力（Self-Attention）。对于给定输入序列，模型通过查询向量（Query）、键向量（Key）和值向量（Value）之间的相似性计算不同位置之间的依赖关系，从而动态聚合上下文信息。多头注意力机制进一步允许模型从多个子空间同时建模序列关系，使其能够有效处理长距离依赖。对于推荐系统中的序列任务而言，用户历史行为序列天然适合被视作一个 token 序列，因此基于 Transformer 的模型在序列推荐中得到广泛应用 \cite{kang2018sasrec}。

随着生成式推荐的发展，研究者开始尝试将用户交互历史组织成文本或结构化 token 序列，并利用大语言模型通过自回归方式直接生成目标物品的表示 \cite{rajput2023tiger,li2024llm4poi,wang2025gnprsid}。在该范式下，推荐任务不再被视作独立的分类或排序问题，而是被重新解释为“给定上下文，生成下一个目标 token 序列”的序列建模问题。这种转换使推荐模型能够更自然地利用大语言模型擅长的上下文推理能力，也为引入离散语义 ID 提供了合理的模型接口。

然而，大语言模型通常参数规模庞大，直接对全部参数进行全量微调成本较高，不仅需要较大的显存和算力支持，也容易造成训练不稳定或过拟合。因此，参数高效微调（Parameter-Efficient Fine-Tuning, PEFT）成为近年来的重要研究方向。PEFT 的基本思想是在尽量冻结预训练模型主体参数的前提下，仅训练少量新增参数或局部参数，从而在显著降低训练成本的同时保持较好的任务适应能力。

LoRA（Low-Rank Adaptation）是参数高效微调中最具代表性的方法之一 \cite {hu2022lora}。该方法认为，在大模型适配下游任务时，权重更新矩阵往往具有低秩特性，因此可以将原有线性变换的参数更新表示为两个低秩矩阵的乘积。设原始权重矩阵为 $W \in \mathbb{R}^{d \times k}$，则 LoRA 不直接更新 $W$，而是引入两个较小矩阵 $A \in \mathbb{R}^{d \times r}$ 和 $B \in \mathbb{R}^{r \times k}$，用它们的乘积近似表示权重增量：
\begin{equation}
W' = W + \Delta W = W + AB,
\end{equation}
其中 $r \ll \min(d,k)$ 为低秩维度。通过这种方式，训练时仅需更新 $A$ 和 $B$，即可在较低参数开销下实现对大模型的有效适配。

% 对于本文研究的 POI 序列预测任务而言，大语言模型的作用主要体现在两个方面。首先，在 GNPR-SID 框架中，经过 RQ-VAE 学习得到的 SID 可以自然表示为离散 token 序列，用户历史访问序列也能够组织为结构化 token 输入，因此大语言模型能够以统一的序列建模方式学习“历史行为--目标 POI 表示”之间的映射关系。其次，参数高效微调使得在有限算力条件下训练大模型成为可能，尤其适合毕业设计或中小规模实验场景。

在本文的方法中，大语言模型并不直接负责学习 POI 的底层语义表示，而主要承担\emph{基于历史访问序列生成目标 SID}的任务。POI 表示学习阶段的重点在于通过 geo\_emb 和 RQ-VAE 获得更合理的离散语义编码；大语言模型阶段的重点则在于利用用户历史序列、时间上下文以及 SID 结构信息完成目标 POI 预测。这种分阶段设计有助于提升模型设计的清晰性与可控性。

\section{地理空间散列编码}
\label{sec:geohash}

地理空间信息是 POI 推荐任务的重要组成部分。相比一般推荐对象，POI 天然携带精确的空间坐标，因此其表示不仅需要反映语义属性，还需要体现地理位置及空间邻近关系。最直接的方式是使用经纬度坐标作为数值特征，但原始经纬度具有连续、高精度和缺乏离散结构等特点，直接用于序列建模或离散 token 建模时往往并不理想。因此，如何将连续地理坐标转换为适合推荐模型使用的离散空间表示，成为地理信息建模中的关键问题。

地理空间散列编码（Geospatial Hash Encoding）是一类常见的空间离散化方法，其代表是 Geohash。Geohash 的基本思想是通过对地理空间进行递归划分，将二维经纬度坐标映射为一个字符串编码。随着编码长度的增加，Geohash 表示的空间区域会逐渐缩小，因此更长的编码意味着更高的位置分辨率。对于任意一个地理坐标 $(lat, lon)$，可以通过 Geohash 编码函数得到长度为 $L$ 的字符串：
\begin{equation}
g = \mathrm{Geohash}(lat, lon, L).
\end{equation}

Geohash 编码的一个重要特性在于其层级性和前缀共享性质。若两个位置在地理空间中相互接近，则它们的 Geohash 编码往往共享更长的公共前缀；反之，若两个位置相距较远，则其前缀通常较早出现差异。这意味着 Geohash 不仅是一种离散编码方式，还天然包含一定的空间邻近结构。对于推荐任务而言，这种特性非常有价值，因为它使模型能够通过共享的编码前缀感知区域相似性和局部空间模式。

在实际应用中，Geohash 常被用作地理位置离散化、空间索引和区域聚合工具。对于本文任务而言，地理空间散列编码具有两方面意义。首先，它能够将连续的经纬度信息转化为可压缩、可比较的离散地理区域表示，从而降低直接处理连续坐标的复杂度。其次，Geohash 编码所蕴含的层级结构可以为后续地理信息增强表示构造提供启发，即通过保留区域相邻性和空间层次性，提高模型对 POI 空间分布模式的表达能力。

% 本文当前方法并不将 GID 作为显式输出目标与 SID 联合生成，而是更强调在表示学习阶段引入地理信息。因此，Geohash 及其他地理空间散列编码方法在本文中的意义主要体现在两个层面：一是为地理信息离散化提供理论依据，说明空间信息可以通过结构化方式嵌入模型；二是为后续地理增强表示设计提供参考，使地理属性能够在 POI 表示学习过程中以更自然的方式融入离散语义编码。

% 从更一般的角度看，地理空间散列编码也可以被理解为一种“空间区域化建模”的方式。它并不追求精确还原绝对坐标，而是更关注不同位置在离散区域空间中的相对关系。对于用户移动行为建模而言，这种区域化思想具有较强合理性。因为用户访问行为通常呈现显著的局部性和区域性，其下一次访问地点更可能位于若干相邻区域之内，而非任意城市位置。因此，将地理空间关系转化为离散区域关系，有助于模型在较低复杂度下捕获用户空间迁移规律。

在本文研究中，虽然核心方法是利用 geo\_emb 替代原始 \texttt{Catname} 和 \texttt{Region} 表示并输入 RQ-VAE，但地理空间散列编码仍然为本文提供了重要的方法启发，即：地理信息应当以结构化方式参与表示学习，而非仅仅作为一个附加数值特征被动输入模型。正是在这一思想下，本文尝试将空间信息前置融入 POI 表示构造过程，使其在后续 SID 学习中发挥更直接的作用。

\section{本章小结}
\label{sec:theory_summary}

本章从四个方面介绍了本文研究所依赖的理论基础与关键技术。首先，对 POI 序列预测任务进行了概述，分析了其相较于一般序列推荐任务所具有的时空约束性、语义复杂性和数据稀疏性特点。其次，介绍了 RQ-VAE 的基本原理，说明其如何通过多级残差量化将连续表示映射为层级化离散标识，并为 SID 学习提供基础。接着，介绍了大语言模型及参数高效微调技术，阐明了其在生成式 POI 序列预测中的作用及采用 LoRA 等方法进行高效适配的必要性。最后，介绍了地理空间散列编码方法，说明了 Geohash 等方法在空间离散表示和区域邻近性建模中的价值。
通过上述分析可以看出，POI 序列预测任务的有效建模不仅依赖于强大的序列生成模型，还依赖于合理的 POI 表示方式。

% 通过上述分析可以看出，POI 序列预测任务的有效建模不仅依赖于强大的序列生成模型，还依赖于合理的 POI 表示方式。尤其是在 GNPR-SID 框架下，若要提升模型对空间迁移规律的感知能力，关键在于使地理信息能够真正参与离散语义表示学习过程。基于此，下一章将进一步介绍本文所使用的数据集、问题定义以及数据预处理流程，并为后续提出地理信息增强的 GNPR-SID 方法做好准备。